\subsection{Why the floor calculations retain their force at later indices}
\label{ext:finite-floor-transfer}
The reuse of a finite resource calculation needs a statement about each
analytic bound, rather than monotonicity of the optimal constants $C_n$.
Let $N$ be its fixed evaluation index and $n\ge N$ the actual index.
Fix one contrary maximizing marginal and all of its acquired moments,
endpoints and signed phase constraints. Throughout a finite refinement
sequence retain
\begin{equation}\label{ext:floor-invariant}
 \mathcal R_n\le r,\qquad nD\le Nd.
\end{equation}
Here $\mathcal R_n$ is the complete raw Fourier remainder, not one
selected summand. The smaller inequality $D\le d$ is available for
index-free moment and gap estimates, but does not replace the second
inequality in \eqref{ext:floor-invariant}.

The index-free operations are intersections of necessary moment
intervals, signwise Cauchy bounds for $\mu=\E X_+$, central-gap
bounds, and physical endpoint prices. For example,
\[
 \mu^2\le p_u(1+\beta_u)/2,\qquad
 \mu^2\le(1-p_l)(1-\beta_l)/2,\qquad
 \mu^2\le\max_{p\in[p_l,p_u]}p(1-p)
\]
do not depend on $n$. If a proved gap inequality gives $6\mu AB\ge P>0$,
then $(A+B)^2\ge2P/(3\mu_u)$, by $4AB\le(A+B)^2$.
Such an improvement requires the already acquired gap and positive mean.
It cannot be used to manufacture a missing initial gap.

For the six-piece high-frequency window bound of Lemma~\ref{ext:high-window50}, keep its admissible
event parameters and window fixed. Its nonconstant index factors are
positive sums and products of
\[
 \sqrt n e^{-cn}\quad(2Nc>1),\qquad n^{-1},\qquad
 e^{-knb^2}/\sqrt n,
\]
and of the decreasing function
\[
 r(x)=\frac{x+2}{x^2+4x+2},\qquad
 r'(x)=-\frac{x^2+4x+6}{(x^2+4x+2)^2}<0.
\]
Thus the ordinary first piece, peak-core correction, peak tails and
zero-peak tail decrease. The normal comparator has the form
$\sqrt n e^{-n/8}\{2r(n/8)+g_c/3\}$, and decreases for $N\ge5$.
The improved first piece has an additional factor $r(n\ell/4)$
with $\ell>0$, and uses its own positive damping and $2Nc>1$ gates.
For the optional improved zero piece, the derivative sign of
$f(x)=\sqrt{x}e^{-x}r(x)$ is the sign of
\[
 4-10x-22x^2-13x^3-2x^4.
\]
This is $-11/128$ at $x=1/4$ and is strictly decreasing for $x>0$.
Hence the optional bound decreases whenever the native gate
$Nkb^2\ge1/4$ holds. All these statements retain the original
window conditions $0<a\le1$, $m_e>0$, $0<d_0\le12$,
$1\le M\le111/10$, $0<w\le1/4$, and positive sine slope.
A minimum of complete valid window bounds stays valid; one need not
assume that the minimizing window remains the same at every index.
Rational numerical upper bounds are taken at the fixed floor $N$.

The low-frequency envelopes have index factors
\[
 (n/(n-1))^{3/2},\quad (n/(n-1))^{3/2}/(n-1),\quad
 a_n^{-j},\quad b_n^{-j},\quad n^{-1/2},
\]
where $a_n=(n-1)/(2n)$, $b_n=(45-M)(1-2/n)/96$ and $j>0$.
The first two decrease and $a_n,b_n$ increase positively on the
retained moment domain. Sums and products with nonnegative fixed
numerators therefore decrease. This applies to the full remainders
$e/(2a_n^2)+B^2/(2b_n^3)$ and $L_2/\sqrt n$.
Every paired-moment eligibility condition remains necessary; the
index comparison cannot repair a failed moment-rectangle hypothesis.

The phase estimate supplies
\[
 \phi(H/\sqrt n)\in[\phi_{0,l}(1-5r/(2N)),2/5]
\]
under its raw-remainder and positivity gates. The actual lower
endpoint increases as $n$ increases. Similarly, for $r,D_l\ge0$,
\[
 r/n-\phi_{0,l}(1-5r/(2n))D_l/6
\]
increases with $1/n$, so evaluation at $N$ is an upper bound.
Put $g_0=(79/25)\rho-3-d$ and $t=(d-4/25)_+$.
The general phase curvature has derivative $(g_0+t)/2\ge0$
with respect to the lower density: if $d\le4/25$, use
$g_0\ge4/25-d$; otherwise $g_0+t=(79/25)(\rho-1)$.
The signed paired curvature is
\[
 C\rho+\tfrac12\min\{\phi_lg_0,(2/5)g_0\}.
\]
It increases with $\phi_l$ for $g_0\ge0$ and is constant in
$\phi_l$ for $g_0<0$. This proves the required direction with
both signs retained. The positive additions and all native
curvature gates are checked before using the phase maximum.

For the sharp contrary score, $C>\CE$ and $\CE<40974/10^5$ give
\[
 n(C-\CE)\rho
 \ge N\max\{0,C-40974/10^5\}\rho_l.
\]
The breakpoint $C=40974/10^5$ is therefore included in the score
cover. On either side the charge is affine in $C$; the same coherent
$(C,\rho)$ corner argument bounds the remaining phase expression.
A charged phase budget $B$ with $nD\le B/\kappa$,
$\kappa=997/15000$, exports $d=B/(\kappa N)$.
A negative budget excludes the law before division. Maximum over a
complete cover, upward rounding and minimum with a previous valid
entry all preserve \eqref{ext:floor-invariant}. Induction over the
finite saved passes proves the invariant without a limiting argument.
A positive $K/n$ calculation is separate: its $K\rho$ charge must
be retained and is not covered by this sharp-target simplification.

Finally, if the physical affine resource is
\[
 D\ge b+g d_{\rm in}+(N/n)h_N\lambda_n,
 \qquad \lambda_n=ns,\quad b,g,h_N\ge0,
\]
then \eqref{ext:floor-invariant} gives
\[
 (n/N)(b+g d_{\rm in})+h_N\lambda_n\le d,
 \qquad b+g d_{\rm in}+h_N\lambda_n\le d.
\]
The original joint intercepts therefore bound the actual noise.
An equivalent directed contract records $D\le Nd/n$, noise price
$Nh_N/n$, and actual phase normalization $15000/(997n)$.
The actual index remains in every binomial mass and count probability.
The present applications are $N=25$, $26\le n\le30$, and $N=31$,
$32\le n\le34$, with a separate support calculation for every
actual index and every support state used by the refinements.

\paragraph{Extracting a phase bound from the whole charged envelope.}
The two local B1 refinements at floor $31$ also tighten the signed
phase intervals. Suppose that, for either sign and $h=|H|\ge0$,
the same acquired law satisfies
\[
 0\le D\le c_n\max_j\{A_j+F_jh-Q_jh^2\},\qquad
 c_n>0,\quad Q_j>0\quad\hbox{for every }j.
\]
The maximum is over the complete charged corner family; its members
cannot be selected independently of the law. At least one corner
polynomial is nonnegative at the actual $h$. Its discriminant
$\Delta_j=F_j^2+4Q_jA_j$ is therefore nonnegative, and solving that
quadratic inequality gives
\[
 h\le\frac{F_j+\sqrt{\Delta_j}}{2Q_j}.
\]
Choose a nonnegative rational enclosure $u_j$ with
$u_j^2\ge\Delta_j$ whenever $\Delta_j\ge0$, and set
\[
 r_j=\begin{cases}
 0,&\Delta_j<0,\\
 \max\{0,(F_j+u_j)/(2Q_j)\},&\Delta_j\ge0.
 \end{cases}
\]
A negative discriminant cannot supply the nonnegative corner, so
$h\le\max_j r_j$. Intersect this bound with the previously acquired
bound for the same sign. The construction is valid at a double root
and at $h=0$. Its checker requires every $Q_j>0$; it has no unchecked
linear or constant branch.

The fixed floor-$31$ normalization is $15000/(997\cdot31)$.
Under the weighted transfer, the actual normalization becomes
$c_n=15000/(997n)$, which stays positive. The implication from
$D\ge0$ to the nonnegativity of at least one charged corner is
therefore unchanged. Separate maxima of the roots of the minus and
plus families give the two signed phase caps used by the actual-index
count contracts.

\paragraph{Transporting the physical variance refinement.}
One resource step uses the bounded conditional supports to obtain
$s\le B\delta$, where $s$ is the rescaled conditional variance,
$\delta=d_{\rm in}$, and
\[
 B=\max(O_0,O_1)\max(p_u,1-p_l)>0.
\]
Here $O_0,O_1$ are outward support offsets in the same conditional
coordinates. If the acquired joint plane is
$\delta/d_0+\lambda_n/L\le1$, with $\lambda_n=ns$, then
\[
 \lambda_n\le \frac{1}{L^{-1}+(nBd_0)^{-1}}.
\]
The physical variance parameter is
$\theta_n=\lambda_n/(n p(1-p)+\lambda_n)$. Writing
\[
 v_l=\min_{p\in[p_l,p_u]}p(1-p)>0
\]
gives
\[
 \theta_n\le
 \frac{1}{1+n v_l/L+v_l/(Bd_0)}
 \le \frac{1}{1+N v_l/L+v_l/(Bd_0)}
 \quad(n\ge N).
\]
The last quantity is precisely the floor calculation's variance
refinement. Downward rounding of the support scale enlarges $O_0,O_1$;
upward rounding of this last bound and intersection with the previous
physical $\theta$ bound preserve its direction. This transports the
bound for $\theta$, rather than treating the floor's diagnostic
$\lambda_N$ as an actual-index bound for $\lambda_n$.
The original joint intercept $L$ remains the one used in the affine
plane. For the selected floor-$25$ first-band resource there are $193$
such saved refinements, each checked at every actual index $26$--$30$.
